\subsection{Finitely Generated PID Modules}

\begin{theorem} \label{thm:finitely-generated-torsion-free-pid-free} \hlt{Finitely Generated Torsion-Free Modules Are Free}\\
For a f.g. $R$-module $M$, $M$ is free $\iff M$ is torsion-free.
\begin{proof}
\begin{enumerate}[label=\roman*.]
\setlength{\itemsep}{0pt}
\item[$(\Rightarrow)$] If $M\cong R^n$, then $am=0$ with $a\neq0$ forces every coordinate of $m$ to be $0$, since $R$ is a domain.
\item[$(\Leftarrow)$] Let $M$ be torsion-free and $G=\{u_1,\ldots,u_n\}$ generate $M$.\\
Choose maximal $R$-independent $S=\{v_1,\ldots,v_k\}\subseteq G$. For $w\in G\setminus S$, dependence of $S\cup\{w\}$ gives $a_w\neq0$ and $b_{w,i}\in R$ with $a_ww+\sum\limits_i b_{w,i}v_i=0$. Put $a=\prod_{w\in G\setminus S}a_w$. Then $aM\subseteq\langle\!\langle S\rangle\!\rangle_R$, so $aM$ is free by Theorem \ref{thm:submodule-of-finite-free-pid-is-free}. Since $M$ is torsion-free, $M\xrightarrow{\sim}aM$, $m\mapsto am$; hence $M$ is free.
\end{enumerate}
\noindent
\hlt{Structural Commentary}: Clear denominators into a free submodule.\\
\hlt{Key Insight}: One scalar embeds $M$ into a free sp	an.
\end{proof}
\end{theorem}

\begin{theorem} \label{thm:finitely-generated-pid-free-plus-torsion} \hlt{Free Plus Torsion Decomposition}\\
For f.g. $M$, $M=M_{\operatorname{tor}}\oplus F$ with $F$ f.g. free. Moreover $M_{\operatorname{tor}}$ is intrinsic and $F\cong M/M_{\operatorname{tor}}$.
\begin{proof} By Proposition \ref{prop:torsion-submodule-over-domain}, $T\defeq M_{\operatorname{tor}}\leq M$ and $Q\defeq M/T$ is f.g. torsion-free.\\
Thus $Q$ is free by Theorem \ref{thm:finitely-generated-torsion-free-pid-free}. Since $Q$ is free, the quotient map $\pi:M\twoheadrightarrow Q$ splits; hence $M=T\oplus\operatorname{im}\sigma$ for any section $\sigma:Q\to M$, and $\operatorname{im}\sigma\cong Q$. If $M=T'\oplus G$ with $T'$ torsion and $G$ free, then $m=t'+g\in M_{\operatorname{tor}}$ forces $g=0$, so $T'=M_{\operatorname{tor}}$. Every complement is isomorphic to $M/M_{\operatorname{tor}}$.\\
\hlt{Structural Commentary}: F.g. PID modules split into torsion plus free.\\
\hlt{Key Insight}: The quotient by torsion is free, and freeness makes the quotient map split.
\end{proof}
\end{theorem}

\begin{definition} \label{def:primary-module-over-pid} \hlt{Primary Module}\\
Let $p\in R$ be prime. A torsion $R$-module $M$ is \hlt{$p$-primary} if $o(M)\sim p^e$ for some $e\geq1$.
\end{definition}

\begin{theorem} \label{thm:primary-decomposition-pid-modules} \hlt{Primary Decomposition}\\
Let $M$ be f.g. torsion with $o(M)=\mu\sim\prod\limits_{i=1}^n p_i^{e_i}$, where the $p_i$ are distinct nonassociate primes.\\
Put $q_i\defeq p_i^{e_i}$, $\mu_i\defeq\mu/q_i$, and $M_i\defeq\mu_iM$. Then $M=\bigoplus\limits_{i=1}^nM_i$, $M_i=\{v\in M:q_iv=0\}$, and $o(M_i)\sim q_i$.\\
The decomposition is unique up to ordering.
\begin{proof}
Since $\gcd(\mu_1,\ldots,\mu_n)\sim1$, choose $a_i\in R$ with $\sum\limits_i a_i\mu_i=1$. Hence $x=\sum\limits_i a_i\mu_ix$, so $M=\sum_iM_i$.\\
Also $q_iM_i=\mu M=0$. Conversely, if $q_iv=0$, choose $a,b\in R$ with $a\mu_i+bq_i=1$; then $v=a\mu_iv\in M_i$.\\
Thus $M_i=\{v:q_iv=0\}$. If $\sum_i x_i=0$ with $x_i\in M_i$, applying $\mu_i$ gives $\mu_ix_i=0$ as $q_j\mid\mu_i$ for $j\neq i$.\\
Since $q_ix_i=0$ and $\gcd(q_i,\mu_i)\sim1$, Bezout gives $x_i=0$; hence the sum is direct.\\
Now $o(M_i)\mid q_i$ and $\operatorname{lcm}_i o(M_i)\sim o(M)\sim\prod_iq_i$, so $o(M_i)\sim q_i$.\\
Finally each summand is the intrinsic set $\{v:q_iv=0\}$, where $q_i$ is the $p_i$-part of $o(M)$.\\
\hlt{Structural Commentary}: Primary decomposition separates a torsion module prime by prime.\\
\hlt{Key Insight}: Bezout idempotents built from the coprime factors of the order project onto the primary pieces.
\end{proof}
\end{theorem}

\begin{lemma} \label{lem:p-torsion-layer-vector-space} \hlt{The First $p$-Torsion Layer}\\
Let $p\in R$ be prime and $M[p]\defeq\ker(p:M\to M)=\{v\in M:pv=0\}$.\\
Then $M[p]$ is an $R/\langle p\rangle$-vector space. If $M=S\oplus T$, then $M[p]=S[p]\oplus T[p]$.
\begin{proof} $\langle p\rangle$ is maximal by Theorem \ref{thm:pid-prime-irreducible-bezout} and Theorem \ref{thm:quotient-by-maximal-ideal-is-field}; hence $\kappa\defeq R/\langle p\rangle$ is a field. Define $(r+\langle p\rangle)v\defeq rv$.\\
If $r-r'\in\langle p\rangle$ and $v\in M[p]$, then $(r-r')v=0$, so the action is well-defined; the axioms are inherited from $M$.\\
If $v=s+t\in M[p]$ under $M=S\oplus T$, then $0=pv=ps+pt$ gives $ps=pt=0$ by directness.\\
Thus $v\in S[p]\oplus T[p]$, and the reverse inclusion is immediate.\\
\hlt{Structural Commentary}: Killing by $p$ turns scalar multiplication modulo $p$ into field scalar multiplication.\\
\hlt{Key Insight}: The layer $M[p]$ counts the number of cyclic primary summands.
\end{proof}
\end{lemma}

\begin{theorem} \label{thm:cyclic-decomposition-primary-pid} \hlt{Cyclic Decomposition of a Primary Module}\\
Let $M$ be f.g. $p$-primary torsion with $o(M)\sim p^e$.\\
Then $M=\bigoplus\limits_{i=1}^m\langle\!\langle v_i\rangle\!\rangle_R$ with $\operatorname{Ann}_R(v_i)=\langle p^{e_i}\rangle$ and $e=e_1\geq e_2\geq\cdots\geq e_m\geq1$.\\
The annihilator chain $\langle p^{e_1}\rangle\subseteq\cdots\subseteq\langle p^{e_m}\rangle$ is intrinsic.
\begin{proof}
\begin{enumerate}[label=(\roman*)]
\setlength{\itemsep}{0pt}
\item \hlt{Existence}: Choose a finite presentation of $M$.\\
Corollary \ref{cor:smith-normal-form-presents-pid-modules} gives $M\cong R^s\oplus\bigoplus\limits_{\ell=1}^tR/\langle d_{\ell}\rangle$, with $d_{\ell}$ nonunit. Since $M$ is torsion, $s=0$.\\
Since $M$ is $p$-primary, $d_{\ell}\sim p^{e_{\ell}}$ with $e_{\ell}\geq1$. Reorder so $e_1\geq\cdots\geq e_m$.\\
Proposition \ref{prop:orders-direct-sum-lcm} gives $o(M)\sim p^{e_1}$, hence $e_1=e$.
\item \hlt{Uniqueness}: Put $\kappa=R/\langle p\rangle$. By Lemma
\ref{lem:p-torsion-layer-vector-space}, each summand $R/\langle p^{e_i}\rangle$ contributes one $\kappa$-dimension to $M[p]$, so $\dim_{\kappa}M[p]=m$. Passing from $M$ to $pM$ replaces each $e_i>1$ by $e_i-1$ and removes each $e_i=1$.\\
Induction on $e$ recovers the exponent multiset, hence the annihilator chain.
\end{enumerate}
\noindent
\hlt{Structural Commentary}: A primary module splits into cyclic strings whose lengths are the exponents $e_i$.\\
\hlt{Key Insight}: The dimensions of the layers $M[p],pM[p],p^2M[p],\ldots$ determine the string lengths.
\end{proof}
\end{theorem}

\begin{proposition} \label{prop:p-primary-layer-counts-recover-exponents} \hlt{Layer Counts Recover Primary Exponents}\\
Let $\kappa=R/\langle p\rangle$ and $\lambda_k(M)\defeq\dim_{\kappa}(p^{k-1}M/p^kM)$ for $k\geq1$.\\
If $M\cong\bigoplus\limits_{i=1}^mR/\langle p^{e_i}\rangle$ with $e_1\geq\cdots\geq e_m\geq1$, then $\lambda_k(M)=\#\{i:e_i\geq k\}$.\\
Hence the number of summands of exact exponent $k$ is $\lambda_k(M)-\lambda_{k+1}(M)$.
\begin{proof} For $C=R/\langle p^e\rangle$, $p^{k-1}C/p^kC$ is $0$ if $k>e$ and is $\kappa$ if $k\leq e$, generated by the class of $p^{k-1}$.\\
Since direct sums commute with $p^{k-1}M$ and $p^kM$, dimensions add, giving $\lambda_k(M)=\#\{i:e_i\geq k\}$.\\
Thus $\lambda_k-\lambda_{k+1}$ counts exactly the summands ending at exponent $k$.\\
\hlt{Structural Commentary}: Successive $p$-layers determine the exponent list.\\
\hlt{Key Insight}: A cyclic summand of order $p^e$ contributes one dimension to exactly the first $e$ layers.
\end{proof}
\end{proposition}
